\documentclass[11pt,a4paper]{article}

% ── Packages ──
\usepackage[utf8]{inputenc}
\usepackage[T1]{fontenc}
\usepackage[english]{babel}
\usepackage[margin=2.2cm]{geometry}
\usepackage{graphicx}
\usepackage{booktabs}
\usepackage{tabularx}
\usepackage{longtable}
\usepackage{array}
\usepackage{xcolor}
\usepackage{colortbl}
\usepackage{hyperref}
\usepackage{enumitem}
\usepackage{fancyhdr}
\usepackage{titlesec}
\usepackage{caption}
\usepackage{float}
\usepackage{multirow}
\usepackage{makecell}
\usepackage{tcolorbox}
\usepackage{pifont}

% ── Colors ──
\definecolor{primary}{HTML}{1A5276}
\definecolor{accent}{HTML}{2E86C1}
\definecolor{lightbg}{HTML}{EBF5FB}
\definecolor{successgreen}{HTML}{27AE60}
\definecolor{warncolor}{HTML}{F39C12}
\definecolor{dangerred}{HTML}{E74C3C}
\definecolor{darktext}{HTML}{2C3E50}

% ── Hyperref ──
\hypersetup{
    colorlinks=true,
    linkcolor=primary,
    urlcolor=accent,
    citecolor=primary,
    pdftitle={Cahier des Charges – SAS-to-Python/PySpark Conversion Accelerator},
    pdfauthor={},
}

% ── Header/Footer ──
\pagestyle{fancy}
\fancyhf{}
\fancyhead[L]{\small\textcolor{primary}{SAS-to-Python/PySpark Conversion Accelerator}}
\fancyhead[R]{\small\textcolor{primary}{Cahier des Charges}}
\fancyfoot[C]{\thepage}
\renewcommand{\headrulewidth}{0.4pt}

% ── Section Formatting ──
\titleformat{\section}{\Large\bfseries\color{primary}}{\thesection}{1em}{}[\vspace{-0.3em}\textcolor{accent}{\rule{\textwidth}{0.8pt}}]
\titleformat{\subsection}{\large\bfseries\color{darktext}}{\thesubsection}{0.8em}{}
\titleformat{\subsubsection}{\normalsize\bfseries\color{accent}}{\thesubsubsection}{0.6em}{}

% ── Custom Commands ──
\newcommand{\cmark}{\textcolor{successgreen}{\ding{51}}}
\newcommand{\xmark}{\textcolor{dangerred}{\ding{55}}}

\newcolumntype{L}[1]{>{\raggedright\arraybackslash}p{#1}}
\newcolumntype{C}[1]{>{\centering\arraybackslash}p{#1}}

% ── tcolorbox Styles ──
\tcbuselibrary{skins,breakable}
\newtcolorbox{infobox}[1][]{
    colback=lightbg, colframe=accent, fonttitle=\bfseries,
    title=#1, boxrule=0.5pt, arc=3pt, left=6pt, right=6pt, top=4pt, bottom=4pt
}
\newtcolorbox{warnbox}[1][]{
    colback=warncolor!5, colframe=warncolor, fonttitle=\bfseries,
    title=#1, boxrule=0.5pt, arc=3pt, left=6pt, right=6pt, top=4pt, bottom=4pt
}

% ══════════════════════════════════════════════════════════════
\begin{document}

% ── TITLE PAGE ──
\begin{titlepage}
\centering
\vspace*{2cm}

{\Huge\bfseries\textcolor{primary}{Cahier des Charges}\\[0.5cm]}
{\LARGE\textcolor{accent}{SAS-to-Python/PySpark Code Conversion Accelerator}\\[1.5cm]}

\textcolor{accent}{\rule{0.6\textwidth}{1.2pt}}\\[1.5cm]

\begin{tabular}{rl}
    \textbf{Date:} & February 2026 \\
\end{tabular}

\vfill
\textcolor{accent}{\rule{0.6\textwidth}{0.4pt}}
\end{titlepage}

% ── TABLE OF CONTENTS ──
\tableofcontents
\newpage

% ══════════════════════════════════════════════════════════════
\section{Context and Problem Statement}
% ══════════════════════════════════════════════════════════════

\subsection{Background}

Many organizations still rely on SAS for their data processing, statistical analysis, and reporting workflows. As the industry moves toward open-source ecosystems (Python, PySpark), there is a growing need to convert legacy SAS codebases into modern Python or PySpark scripts.

Manual conversion is slow, error-prone, and requires deep expertise in \textbf{both} SAS and Python. A single analyst can convert only a few files per day, and subtle semantic differences between the two languages lead to silent bugs that are hard to detect.

\subsection{The Core Challenge}

The challenge is not just about translating syntax — it is about preserving \textbf{semantic equivalence}. SAS and Python handle certain operations in fundamentally different ways, which causes silent errors if not handled carefully:

\begin{table}[H]
\centering
\caption{Key Semantic Differences Between SAS and Python}
\renewcommand{\arraystretch}{1.3}
\begin{tabularx}{\textwidth}{@{} l X @{}}
\toprule
\textbf{Area} & \textbf{Why It Is Difficult} \\
\midrule
Date Arithmetic & SAS uses a different date origin than Python, so all date calculations must be adjusted. \\
Merge / Join Logic & SAS merging follows a sorted-join model that behaves differently from Python's merge functions. \\
Row-to-Row State & SAS can carry values from one row to the next automatically; Python has no built-in equivalent. \\
Group Processing & SAS provides automatic first/last row detection within groups; Python requires explicit logic. \\
Missing Values & SAS treats missing values differently in comparisons; Python's missing values behave differently. \\
Output Structure & Some SAS procedures produce specific output table structures that must be replicated exactly. \\
\bottomrule
\end{tabularx}
\end{table}

\subsection{Proposed Solution}

We propose building an \textbf{AI-powered conversion accelerator} — an intelligent pipeline that takes SAS files as input and produces equivalent Python or PySpark scripts as output. The system combines:

\begin{itemize}[leftmargin=1.5em]
    \item \textbf{Intelligent partitioning} — Breaking SAS files into meaningful, self-contained blocks using a research-backed clustering approach (RAPTOR, ICLR 2024).
    \item \textbf{Three retrieval strategies} — Not all blocks need the same amount of help. Simple blocks use Static RAG (fixed retrieval). Blocks with cross-file dependencies use GraphRAG (the system reads the dependency graph to understand context). Complex blocks use Agentic RAG (the system autonomously decides how much to retrieve, reformulates its queries on failure, and retries retrieval — not just generation).
    \item \textbf{Quality assurance} — Translated code goes through automated validation before being accepted.
    \item \textbf{Structured reassembly} — Translated blocks are merged back into a complete, runnable Python script with proper imports and variable resolution.
    \item \textbf{Continuous learning} — Corrections are fed back into the knowledge base, so the system improves over time.
\end{itemize}

\begin{figure}[H]
    \centering
    \includegraphics[width=\textwidth]{fig_solution_overview.png}
    \caption{Solution Overview — From SAS Input to Python Output}
    \label{fig:solution_overview}
\end{figure}


% ══════════════════════════════════════════════════════════════
\section{Objectives}
% ══════════════════════════════════════════════════════════════

\subsection{Primary Objectives}

\begin{table}[H]
\centering
\caption{Primary Objectives}
\renewcommand{\arraystretch}{1.3}
\begin{tabularx}{\textwidth}{@{} c L{3cm} X @{}}
\toprule
\textbf{ID} & \textbf{Objective} & \textbf{Description} \\
\midrule
O1 & End-to-end automation & Accept any valid SAS file and produce an equivalent, syntactically valid Python or PySpark script without manual intervention for standard constructs. \\
O2 & Semantic correctness & Generate Python code that behaves the same as the original SAS code, with explicit handling of the known semantic differences between the two languages. \\
O3 & Broad coverage & Handle all major SAS categories: data processing, statistics, SQL, reporting, macros, and file operations. \\
O4 & Measurable quality & Every component has clear quality criteria with automated evaluation so we can objectively measure how well the system performs. \\
O5 & Research contribution & Validate the RAPTOR-based partitioning approach through a controlled experiment comparing it against a simpler baseline. \\
\bottomrule
\end{tabularx}
\end{table}

\subsection{Secondary Objectives}

\begin{table}[H]
\centering
\caption{Secondary Objectives}
\renewcommand{\arraystretch}{1.3}
\begin{tabularx}{\textwidth}{@{} c L{3.5cm} X @{}}
\toprule
\textbf{ID} & \textbf{Objective} & \textbf{Description} \\
\midrule
O6 & Continuous learning & The system improves over time by incorporating corrections back into its knowledge base automatically. \\
O7 & Fault tolerance & If the system crashes during a large project, it can resume from where it stopped instead of restarting from scratch. \\
O8 & Fully open-source & The entire stack uses free, open-source tools — no proprietary dependencies, deployable at zero cost. \\
\bottomrule
\end{tabularx}
\end{table}


% ══════════════════════════════════════════════════════════════
\section{System Architecture}
% ══════════════════════════════════════════════════════════════

The system is organized as a \textbf{six-phase pipeline}, where each phase transforms the SAS input progressively until a complete Python script is produced. A continuous learning loop feeds back into the knowledge base.

\begin{figure}[H]
    \centering
    \includegraphics[width=\textwidth]{fig_pipeline.png}
    \caption{End-to-End Processing Pipeline}
    \label{fig:pipeline}
\end{figure}

\subsection{Architecture Layers}

The pipeline is implemented across several layers, each with a clear responsibility:

\begin{table}[H]
\centering
\caption{System Layers Overview}
\renewcommand{\arraystretch}{1.35}
\small
\begin{tabularx}{\textwidth}{@{} c l X @{}}
\toprule
\textbf{Layer} & \textbf{Name} & \textbf{Responsibility} \\
\midrule
L2-A & Entry \& Scan & Discover SAS files, detect encoding, register metadata, resolve cross-file dependencies, extract data lineage. \\
L2-B & Streaming Core & Read files asynchronously and parse SAS code line by line, tracking the current state (block type, nesting, macros). \\
L2-C & RAPTOR Partitioning & Detect block boundaries, embed blocks into a vector space, cluster related blocks together, build a hierarchical tree. \\
L2-D & Complexity \& Strategy & Analyze each block's complexity and decide which translation strategy and model to use. \\
L2-E & Persistence \& Index & Store all intermediate results, build a dependency graph, and create checkpoints for crash recovery. \\
L3 & Translation & Convert each SAS block to Python or PySpark, guided by the knowledge base, with cross-verification. \\
L4 & Merge & Reassemble translated blocks into a complete script with proper imports, variable names, and ordering. \\
CL & Continuous Learning & Ingest corrections, monitor quality trends, and trigger knowledge base updates when needed. \\
\bottomrule
\end{tabularx}
\end{table}

\begin{figure}[H]
    \centering
    \includegraphics[width=\textwidth]{fig_architecture_layers.png}
    \caption{Architecture Layers and Data Flow}
    \label{fig:architecture_layers}
\end{figure}


% ══════════════════════════════════════════════════════════════
\section{Component Descriptions}
% ══════════════════════════════════════════════════════════════

\subsection{Phase 1 — File Ingestion (L2-A)}

The first phase discovers and registers all SAS files in the project:

\begin{itemize}[leftmargin=1.5em]
    \item \textbf{File Discovery} — Recursively scans the project folder for all \texttt{.sas} files.
    \item \textbf{Encoding Detection} — Automatically detects the character encoding of each file to avoid garbled text.
    \item \textbf{Pre-Validation} — Checks for structural issues (e.g., unclosed macros, truncated procedures) before processing.
    \item \textbf{Cross-File Dependencies} — Detects when one SAS file references another (via \texttt{\%INCLUDE}, \texttt{LIBNAME}, etc.) and records these relationships.
    \item \textbf{Data Lineage} — Tracks which datasets each file reads from and writes to, building a table-level data flow map.
    \item \textbf{Deduplication} — Uses content hashing to ensure the same file is never processed twice, even after a restart.
\end{itemize}

\subsection{Phase 2 — Code Analysis (L2-B)}

Each SAS file is read and analyzed line by line using an asynchronous streaming approach:

\begin{itemize}[leftmargin=1.5em]
    \item \textbf{Async Streaming} — Files are read in small chunks to keep memory usage low, even for very large files.
    \item \textbf{Stateful Parsing} — A finite state machine (FSM) tracks the current block type (DATA step, PROC, macro, etc.), nesting depth, macro stack, and variable scope as each line is processed.
    \item \textbf{Macro Integrity} — The parser maintains a dedicated macro stack to ensure macros are never accidentally split. A macro block only closes when the matching \texttt{\%MEND} is found — inner \texttt{RUN;} or \texttt{QUIT;} statements do not break the macro boundary.
\end{itemize}

\begin{figure}[H]
    \centering
    \includegraphics[width=\textwidth]{fig_macro_handling.png}
    \caption{Macro Integrity — How Macros Are Kept Together}
    \label{fig:macro_handling}
\end{figure}

\subsection{Phase 3 — Intelligent Partitioning (L2-C)}

This is the core innovation of the system. Instead of splitting code with simple keyword matching, we use a research-backed hierarchical clustering approach:

\begin{itemize}[leftmargin=1.5em]
    \item \textbf{Boundary Detection} — A grammar-based parser identifies where each SAS block starts and ends. For ambiguous cases, a language model resolves the boundary.
    \item \textbf{Block Typing} — Each block is classified into one of nine canonical types: DATA step, PROC block, macro definition, macro invocation, SQL block, conditional block, loop block, global statement, or include reference.
    \item \textbf{Semantic Embedding} — Each block is converted into a numerical vector that captures its meaning, not just its text.
    \item \textbf{Hierarchical Clustering (RAPTOR)} — Related blocks are grouped together using Gaussian Mixture Models. Groups are then summarized and re-clustered recursively, building a tree structure from individual blocks (leaves) up to file-level summaries (root).
\end{itemize}

\begin{figure}[H]
    \centering
    \includegraphics[width=\textwidth]{fig_raptor_tree.png}
    \caption{RAPTOR Hierarchical Clustering — From Blocks to Tree}
    \label{fig:raptor_tree}
\end{figure}

\subsection{Phase 3b — Complexity Analysis \& Strategy Selection (L2-D)}

Each partition is analyzed for its complexity, and a translation strategy is selected:

\begin{itemize}[leftmargin=1.5em]
    \item \textbf{Complexity Features} — Multiple features are extracted: cyclomatic complexity, nesting depth, macro density, cross-file dependencies, and SQL complexity.
    \item \textbf{Risk Classification} — A calibrated machine learning model assigns each block a risk level (low, moderate, or high). The confidence of this prediction is also computed.
    \item \textbf{Model Routing} — Low-risk blocks are sent to a smaller, faster language model. Moderate and high-risk blocks are sent to a larger, more capable model. Blocks where the system is uncertain are flagged for human review.
\end{itemize}

\begin{figure}[H]
    \centering
    \includegraphics[width=\textwidth]{fig_risk_routing.png}
    \caption{Risk-Based Model Routing}
    \label{fig:risk_routing}
\end{figure}

\subsection{Persistence \& Dependency Graph (L2-E)}

All intermediate results are stored persistently so the pipeline can resume after interruptions:

\begin{itemize}[leftmargin=1.5em]
    \item \textbf{Structured Storage} — Partition data, conversion results, and metadata are stored in a relational database.
    \item \textbf{Vector Index} — Block embeddings and RAPTOR tree nodes are stored in a vector database for fast semantic retrieval during translation.
    \item \textbf{Dependency Graph} — Cross-file and cross-block dependencies are stored in a graph database. Circular dependencies are detected and handled by grouping them into translation units.
    \item \textbf{Checkpoints} — The system saves progress at regular intervals so it can resume from the last checkpoint after a crash.
    \item \textbf{Audit Logging} — Every operation is logged for traceability and debugging.
\end{itemize}

\subsection{Phase 4 — Translation (L3)}

Each SAS partition is translated using one of three retrieval-augmented generation (RAG) paradigms, selected automatically based on the block's complexity and dependencies:

\subsubsection*{Static RAG — Simple Blocks}
\begin{itemize}[leftmargin=1.5em]
    \item Used for \textbf{LOW risk} partitions with no cross-file dependencies.
    \item Fixed retrieval: 3 similar SAS-to-Python pairs from the knowledge base (leaf level).
    \item Single-pass generation using a small, fast model. No cross-verification.
\end{itemize}

\subsubsection*{GraphRAG — Dependency-Aware Blocks}
\begin{itemize}[leftmargin=1.5em]
    \item Used when the block has \textbf{cross-file dependencies}, is part of an SCC (circular dependency group), or uses macros defined elsewhere.
    \item Before querying the knowledge base, the system traverses the dependency graph to find upstream partitions that have already been translated.
    \item These upstream translations are injected as additional context — the language model sees what its dependencies produced.
    \item Retrieval from knowledge base at \textbf{cluster level} (related procedure families), providing broader context.
\end{itemize}

\subsubsection*{Agentic RAG — Complex / High-Risk Blocks}
\begin{itemize}[leftmargin=1.5em]
    \item Used for \textbf{MODERATE or HIGH risk} blocks, or when a failure mode is detected.
    \item The agent \textbf{autonomously controls retrieval}: it decides whether to retrieve, how many results to fetch (adaptive $k$: 5\textrightarrow8\textrightarrow11), and which level of the knowledge base hierarchy to query.
    \item If cross-verification fails, the agent \textbf{retries retrieval} — not just generation. It reformulates the query (adding the specific error and failure mode), increases $k$, and escalates to a higher level in the RAPTOR hierarchy.
    \item Final escalation: root-level retrieval combined with graph context from upstream dependencies.
    \item If all retries fail, the block is marked as needing human review.
\end{itemize}

\begin{figure}[H]
    \centering
    \includegraphics[width=\textwidth]{fig_rag_paradigms.png}
    \caption{Three RAG Paradigms — Static, GraphRAG, and Agentic RAG}
    \label{fig:rag_paradigms}
\end{figure}

After translation, \textbf{all paradigms} pass through the same quality gate:
\begin{itemize}[leftmargin=1.5em]
    \item \textbf{Validation} — The generated code is executed on synthetic data to catch runtime errors.
\end{itemize}

\subsection{Phase 5 — Merge \& Output (L4)}

Translated blocks are reassembled into a complete Python script:

\begin{itemize}[leftmargin=1.5em]
    \item \textbf{Import Consolidation} — All imports are collected, deduplicated, and organized following Python conventions.
    \item \textbf{Variable Resolution} — SAS dataset names are mapped to consistent Python variable names across all blocks.
    \item \textbf{Block Assembly} — Translated blocks are placed in the correct execution order.
    \item \textbf{Syntax Validation} — The final script is checked to ensure it is syntactically valid Python.
    \item \textbf{Report Generation} — A detailed conversion report is generated showing what was converted, what needs review, and quality metrics.
\end{itemize}

\subsection{Phase 6 — Continuous Learning (CL)}

The system improves over time through a feedback loop:

\begin{itemize}[leftmargin=1.5em]
    \item \textbf{Correction Ingestion} — When a translation is corrected (by a user or automatically), the correction is verified and added to the knowledge base.
    \item \textbf{Quality Monitoring} — The system tracks conversion quality over time and raises alerts if quality drops in any area.
    \item \textbf{Retraining Triggers} — When enough new corrections accumulate or quality metrics drift, the system triggers an update to its models and knowledge base.
\end{itemize}

\begin{figure}[H]
    \centering
    \includegraphics[width=\textwidth]{fig_continuous_learning.png}
    \caption{Continuous Learning Feedback Loop}
    \label{fig:continuous_learning}
\end{figure}


% ══════════════════════════════════════════════════════════════
\section{Knowledge Base}
% ══════════════════════════════════════════════════════════════

The knowledge base is the \textbf{primary quality driver} of the system. Every translation is augmented with verified SAS-to-Python pairs retrieved from this curated collection.

\subsection{How It Was Built}

The knowledge base was constructed \textbf{entirely without a SAS runtime}. A multi-step language model pipeline generates SAS-to-Python pairs, and a separate verification step cross-checks each pair for correctness. Only pairs that pass verification are accepted.

\subsection{What It Covers}

The knowledge base covers all major SAS construct categories:

\begin{table}[H]
\centering
\caption{Knowledge Base Coverage}
\renewcommand{\arraystretch}{1.3}
\begin{tabularx}{\textwidth}{@{} l X @{}}
\toprule
\textbf{Category} & \textbf{What It Includes} \\
\midrule
Data Processing & Basic data steps, assignments, filtering, formatting, arrays \\
Merge \& Joins & All merge types (one-to-one, one-to-many), update operations \\
Row-to-Row State & RETAIN statements, running totals, lag patterns \\
Group Processing & BY-group processing with first/last row detection \\
Date Arithmetic & Date creation, intervals, differences, date parts \\
SQL Operations & SELECT, JOIN, subqueries, aggregations \\
Statistical Procedures & Means, frequencies, regression, logistic models \\
Macros & Definitions, parameters, conditional macros, nested macros \\
File I/O & Import/export for CSV, Excel, and other formats \\
Missing Values & Missing value detection, comparisons, handling patterns \\
\bottomrule
\end{tabularx}
\end{table}

Additionally, the knowledge base includes a set of \textbf{targeted failure-mode examples} — pairs specifically designed to teach the language model how to handle each of the six semantic pitfalls correctly.

\begin{figure}[H]
    \centering
    \includegraphics[width=\textwidth]{fig_knowledge_base.png}
    \caption{Knowledge Base Structure and Retrieval Flow}
    \label{fig:knowledge_base}
\end{figure}

\subsection{Versioning and Evolution}

Every knowledge base entry carries a version number. When a correction is made, the old entry is kept (for traceability) and a new, corrected version is created. A changelog records every modification with timestamps and authorship, and any entry can be rolled back to a previous version if needed.


% ══════════════════════════════════════════════════════════════
\section{Agent Architecture}
% ══════════════════════════════════════════════════════════════

The system is built as a \textbf{multi-agent architecture} where each agent has a single, well-defined responsibility. This design ensures that each component can be developed, tested, and improved independently.

\begin{table}[H]
\centering
\caption{System Agents — One Responsibility Each}
\renewcommand{\arraystretch}{1.25}
\small
\begin{tabularx}{\textwidth}{@{} c l l X @{}}
\toprule
\textbf{\#} & \textbf{Agent} & \textbf{Layer} & \textbf{What It Does} \\
\midrule
1 & FileAnalysisAgent & Entry & Discovers files, detects encoding, runs pre-validation \\
2 & CrossFileDependencyResolver & Entry & Finds cross-file references (\texttt{\%INCLUDE}, \texttt{LIBNAME}) \\
3 & RegistryWriterAgent & Entry & Writes file metadata to the registry, prevents duplicates \\
4 & DataLineageExtractor & Entry & Maps which datasets are read and written \\
5 & StreamAgent & Streaming & Reads files asynchronously in small chunks \\
6 & StateAgent & Streaming & Tracks parsing state (block type, nesting, macros) \\
7 & BoundaryDetectorAgent & Partitioning & Detects where blocks start and end \\
8 & RAPTORPartitionAgent & Partitioning & Clusters blocks and builds the hierarchical tree \\
9 & ComplexityAgent & Strategy & Extracts complexity features from each block \\
10 & StrategyAgent & Strategy & Assigns risk level and selects translation strategy \\
11 & PersistenceAgent & Storage & Saves intermediate results to the database \\
12 & IndexAgent & Storage & Builds and maintains the dependency graph \\
13 & TranslationAgent & Translation & Converts SAS to Python/PySpark using the KB \\
14 & ValidationAgent & Translation & Tests generated code on synthetic data \\
15 & ReportAgent & Output & Generates the conversion report \\
16 & PartitionOrchestrator & Pipeline & Coordinates all agents, manages checkpoints \\
\bottomrule
\end{tabularx}
\end{table}

\begin{figure}[H]
    \centering
    \includegraphics[width=\textwidth]{fig_agent_architecture.png}
    \caption{Agent Architecture — 16 Agents Across All Layers}
    \label{fig:agent_architecture}
\end{figure}

\subsection{Fault Tolerance}

Every agent that depends on an external service (language model API, database, etc.) has a built-in retry mechanism with fallback behavior. If retries are exhausted, the agent gracefully degrades — for example, falling back to a simpler algorithm or flagging the block for human review — rather than crashing the entire pipeline.


% ══════════════════════════════════════════════════════════════
\section{Technology Stack}
% ══════════════════════════════════════════════════════════════

All technologies are \textbf{free and open-source}. The stack was selected for reliability, simplicity, and zero-cost deployment:

\begin{table}[H]
\centering
\caption{Technology Stack}
\renewcommand{\arraystretch}{1.3}
\small
\begin{tabularx}{\textwidth}{@{} l l X @{}}
\toprule
\textbf{Role} & \textbf{Technology} & \textbf{Why This Choice} \\
\midrule
Language Model (light) & Llama 8B / Ollama & Free, runs locally, good enough for simple blocks \\
Language Model (heavy) & Llama 70B / Groq & Free tier API, much more capable for complex blocks \\
Embedding Model & Nomic Embed v1.5 & High quality, runs locally, open license \\
Vector Database & LanceDB & Persistent, supports filtered search, fully embedded \\
Graph Library & NetworkX & Pure Python, in-memory, built-in SCC and topological sort \\
SAS Parser & Lark (LALR) & Handles nested structures, pure Python \\
Complexity Analysis & Radon & Standard complexity metrics out of the box \\
ML Classification & scikit-learn & Simple, interpretable, well-calibrated models \\
Pipeline Orchestration & LangGraph & Built-in state management, branching, and checkpoints \\
Caching / Checkpoints & Redis & Fast, atomic operations, automatic expiration \\
Analytics \& Audit & DuckDB & Fast analytical queries on logs and metrics \\
Validation & Pydantic v2 & Robust data validation, industry standard \\
\bottomrule
\end{tabularx}
\end{table}


% ══════════════════════════════════════════════════════════════
\section{Evaluation Approach}
% ══════════════════════════════════════════════════════════════

\subsection{Code Quality Metric — CodeBLEU}

To evaluate the quality of generated code, we use \textbf{CodeBLEU} (Ren et al., 2020), a metric specifically designed for code generation. Unlike standard text similarity metrics, CodeBLEU considers:

\begin{itemize}[leftmargin=1.5em]
    \item \textbf{Token overlap} — How similar the generated tokens are to the reference.
    \item \textbf{Keyword awareness} — Python keywords are weighted more heavily than variable names.
    \item \textbf{Syntax structure} — Whether the generated code has the same syntactic structure as the reference.
    \item \textbf{Data flow} — Whether the data flows through the code in the same way.
\end{itemize}

\subsection{Evaluation Per Layer}

Each layer of the system has its own set of quality criteria:

\begin{table}[H]
\centering
\caption{Quality Criteria by Layer}
\renewcommand{\arraystretch}{1.35}
\begin{tabularx}{\textwidth}{@{} l X @{}}
\toprule
\textbf{Layer} & \textbf{What We Measure} \\
\midrule
Entry \& Scan & File discovery completeness, encoding accuracy, cross-file resolution coverage, pre-validation accuracy \\
Streaming Core & Processing throughput, memory efficiency, parsing accuracy for block types and nesting \\
RAPTOR Partitioning & Boundary detection accuracy, partition type correctness, retrieval quality (RAPTOR vs. flat baseline) \\
Complexity \& Strategy & Calibration quality, risk classification accuracy, strategy selection correctness \\
Persistence & Deduplication correctness, dependency graph completeness, checkpoint reliability \\
Translation & CodeBLEU score, translation success rate, validation pass rate, human review rate \\
Merge & Syntax validity, import correctness, variable consistency, block ordering \\
Continuous Learning & Knowledge base growth rate, quality improvement over time \\
\bottomrule
\end{tabularx}
\end{table}

\subsection{Ablation Study}

To validate the value of the RAPTOR-based approach, we run a controlled experiment with three configurations:

\begin{enumerate}[leftmargin=1.5em]
    \item \textbf{Baseline} — Language model alone, no knowledge base, no failure-mode detection.
    \item \textbf{KB only} — Knowledge base retrieval added, but no failure-mode detection.
    \item \textbf{Full system} — Knowledge base + failure-mode detection + RAPTOR partitioning.
\end{enumerate}

This allows us to measure exactly how much each component contributes to translation quality.


% ══════════════════════════════════════════════════════════════
\section{Data Storage}
% ══════════════════════════════════════════════════════════════

The system uses multiple specialized databases, each chosen for a specific type of data:

\begin{table}[H]
\centering
\caption{Database Architecture}
\renewcommand{\arraystretch}{1.3}
\begin{tabularx}{\textwidth}{@{} l X @{}}
\toprule
\textbf{Database} & \textbf{What It Stores} \\
\midrule
SQLite & Core data: file registry, partitions, conversion results, cross-file dependencies, data lineage \\
DuckDB & Audit logs: every LLM call, calibration events, quality metrics, feedback history \\
LanceDB & Vector data: block embeddings, RAPTOR tree nodes, knowledge base entries \\
NetworkX & Graph data: partition dependencies, macro call relationships, cross-file edges, SCC groups \\
Redis & Temporary checkpoints for crash recovery \\
\bottomrule
\end{tabularx}
\end{table}

\begin{figure}[H]
    \centering
    \includegraphics[width=\textwidth]{fig_database_architecture.png}
    \caption{Data Storage — Four Databases + Graph Library}
    \label{fig:database_architecture}
\end{figure}


% ══════════════════════════════════════════════════════════════
\section{Timeline}
% ══════════════════════════════════════════════════════════════

The project spans 14 weeks, organized into four priority tiers:

\begin{table}[H]
\centering
\caption{Project Timeline — 14 Weeks}
\renewcommand{\arraystretch}{1.3}
\small
\begin{tabularx}{\textwidth}{@{} c X c @{}}
\toprule
\textbf{Weeks} & \textbf{Deliverable} & \textbf{Priority} \\
\midrule
\rowcolor{successgreen!8} 1–2 & File ingestion, cross-file dependencies, data lineage, gold standard corpus & P0 \\
\rowcolor{successgreen!8} 2–3 & Streaming core: async reader + stateful parser & P0 \\
\rowcolor{successgreen!8} 3–4 & Boundary detection with grammar parser + LLM fallback & P0 \\
\rowcolor{successgreen!8} 4 & Complexity analysis + strategy selection + RAPTOR models & P0 \\
\rowcolor{warncolor!8} 5–6 & Embedding + clustering + tree builder (RAPTOR pipeline) & P1 \\
\rowcolor{warncolor!8} 7 & Persistence layer + dependency graph + audit logging & P1 \\
\rowcolor{warncolor!8} 8 & Pipeline orchestration + checkpoints + fault tolerance & P1 \\
9 & Robustness testing + knowledge base generation (initial set) & P2 \\
10 & Translation layer implementation + KB expansion & P2 \\
11 & Merge layer + continuous learning + KB completion & P2 \\
12 & Ablation study: RAPTOR vs. flat retrieval & P2 \\
\rowcolor{lightbg} 13 & Defense preparation: slides + demo & P3 \\
\rowcolor{lightbg} 14 & Buffer: polish, additional examples, documentation & P3 \\
\bottomrule
\end{tabularx}
\end{table}

\begin{table}[H]
\centering
\caption{Priority Legend}
\renewcommand{\arraystretch}{1.2}
\begin{tabularx}{\textwidth}{@{} l X @{}}
\toprule
\textbf{Priority} & \textbf{Scope} \\
\midrule
\rowcolor{successgreen!8} \textbf{P0} & Core infrastructure — must be completed first. \\
\rowcolor{warncolor!8} \textbf{P1} & Semantic layer + persistence — builds on top of P0. \\
\textbf{P2} & Translation, merge, learning, and evaluation. \\
\rowcolor{lightbg} \textbf{P3} & Defense preparation and polish. \\
\bottomrule
\end{tabularx}
\end{table}


% ══════════════════════════════════════════════════════════════
\section{Risk Management}
% ══════════════════════════════════════════════════════════════

\begin{table}[H]
\centering
\caption{Risk Register}
\renewcommand{\arraystretch}{1.3}
\small
\begin{tabularx}{\textwidth}{@{} L{3.2cm} c c X @{}}
\toprule
\textbf{Risk} & \textbf{Likelihood} & \textbf{Impact} & \textbf{Mitigation} \\
\midrule
API quota exceeded & Medium & High & Multi-tier fallback: cloud API → local large model → heuristic rules. \\
RAPTOR experiment inconclusive & Low & Medium & Report honestly as a negative result. The system still functions without the clustering advantage. \\
Clustering fails on large files & Medium & High & Automatic fallback to simple flat partitioning after repeated failures. \\
Timeline slips & High & Medium & Weekly progress tracking. Scope reduction plan: cut advanced features, never cut core translation. \\
Knowledge base quality issues & Medium & High & Increase verification passes. Human review queue for low-confidence pairs. \\
External tool installation issues & Low & Medium & Every external dependency has a fallback using a simpler alternative. \\
\bottomrule
\end{tabularx}
\end{table}


% ══════════════════════════════════════════════════════════════
\section{Methodology}
% ══════════════════════════════════════════════════════════════

\begin{table}[H]
\centering
\caption{Development Methodology}
\renewcommand{\arraystretch}{1.3}
\begin{tabularx}{\textwidth}{@{} l X @{}}
\toprule
\textbf{Aspect} & \textbf{Approach} \\
\midrule
Process & Hybrid Kanban with weekly sprint reviews. Limited work in progress to stay focused. \\
Version control & Git with clear commit conventions (feature, fix, test, docs). \\
Testing & Each agent has its own unit tests. Integration tests validate the full pipeline. \\
Regression guards & Automated checks ensure that new changes don't break existing functionality. \\
Documentation & Code documentation, architecture diagrams, and a final defense report. \\
\bottomrule
\end{tabularx}
\end{table}


% ══════════════════════════════════════════════════════════════
\section{Deliverables}
% ══════════════════════════════════════════════════════════════

\subsection{Code Deliverables}

\begin{table}[H]
\centering
\caption{Code Deliverables}
\renewcommand{\arraystretch}{1.25}
\begin{tabularx}{\textwidth}{@{} L{4.5cm} X @{}}
\toprule
\textbf{Deliverable} & \textbf{Description} \\
\midrule
Core pipeline package & All agents implementing the six-phase pipeline \\
Test suite & Unit tests for every agent + integration tests for the full pipeline \\
Translation layer & SAS-to-Python/PySpark conversion with KB augmentation \\
Merge pipeline & Import consolidation, variable resolution, script assembly \\
Continuous learning module & Feedback ingestion and quality monitoring \\
\bottomrule
\end{tabularx}
\end{table}

\subsection{Data Deliverables}

\begin{table}[H]
\centering
\caption{Data Deliverables}
\renewcommand{\arraystretch}{1.25}
\begin{tabularx}{\textwidth}{@{} L{4.5cm} X @{}}
\toprule
\textbf{Deliverable} & \textbf{Description} \\
\midrule
Knowledge base & Curated collection of verified SAS-to-Python/PySpark pairs covering all major categories \\
Gold standard corpus & Representative SAS files spanning all complexity levels, used for benchmarking \\
Labeled training data & Training examples for the complexity classifier \\
RAPTOR tree & Hierarchical index built from the gold standard files \\
Dependency graph & Complete cross-file and cross-block dependency graph \\
Ablation results & Experimental comparison: RAPTOR vs. flat retrieval \\
\bottomrule
\end{tabularx}
\end{table}

\subsection{Documentation Deliverables}

\begin{table}[H]
\centering
\caption{Documentation Deliverables}
\renewcommand{\arraystretch}{1.25}
\begin{tabularx}{\textwidth}{@{} L{4.5cm} X @{}}
\toprule
\textbf{Deliverable} & \textbf{Description} \\
\midrule
Architecture documentation & Full system architecture with diagrams and data flow descriptions \\
Conversion reports & Per-file reports showing what was converted, quality metrics, and review flags \\
Defense presentation & Slides + live demo of the full pipeline \\
\bottomrule
\end{tabularx}
\end{table}




\end{document}
